\begin{name}
	{\tenchude}
	{TOÁN 10}
	{LỚP TOÁN THẦY PHÁT}
	{I can accept failure, everyone fails at something. But I can’t accept not trying.}
\end{name}
\TN
\begin{ex}%[0D6N3-3]
Nhiệt độ $\left({}^{\circ}\mathrm{C}\right)$ trung bình các ngày trong tuần đầu tháng $3$ của Nam Định là
\[
\begin{array}{lllllll}
20 & 19 & 22 & 19 & 24 & 28 & 25
\end{array}
\]
Tìm số trung vị của mẫu số liệu trên.
\choice
{\True $22$}
{$23$}
{$24$}
{$25$}
\loigiai{
Sắp xếp theo thứ tự không giảm:
$\begin{array}{ccccccc}19 & 19 & 20 & 22 & 24 & 25 & 28.\end{array}$
Mẫu số liệu trên có $7$ số.\\
Số trung vị là số đứng ở vị trí thứ $\dfrac{7+1}{2}=4$.\\
Vậy số trung vị là $M_e=22$}
\end{ex}

\begin{ex}%[1-HK1-CT-1-BuiThiXuan-HCM-2324]%[VN-MT-6, VM012]%[0D6N4-2]
Một trường THPT đo chiều cao của các bạn học sinh lớp $11$ và thu được mẫu số liệu sau (đơn vị: cm):
\begin{center}
\begin{tabular}{|c|c|c|c|c|c|c|c|}
\hline $161$& $163$ & $175$ & $179$ & $164$ & $170$ & $172$ & $171$ \\
\hline $165$ & $166$ & $166$ & $178$ & $162$ & $160$ & $158$ & $158$ \\
\hline $162$ & $162$ & $174$ & $155$ & $157$ & $163$ & $168$ & 1$64$ \\
\hline
\end{tabular}
\end{center}
Tìm khoảng biến thiên của mẫu số liệu trên.
\choice
{$21$}
{$23$}
{\True $24$}
{$22$}
\loigiai{Khoảng biến thiên của mẫu số liệu trên là $179-155=24$.}
\end{ex}

\begin{ex}%[0D8H1-3]
Bình $A$ chứa $3$ quả cầu xanh, $4$ quả cầu đỏ và $5$ quả cầu trắng. Bình $B$ chứa $4$ quả cầu xanh, $3$ quả cầu đỏ và $6$ quả cầu trắng. Bình $C$ chứa $5$ quả cầu xanh, $5$ quả cầu đỏ và $2$ quả cầu trắng. Từ mỗi bình lấy một quả cầu. Có bao nhiêu cách lấy để cuối cùng được $3$ quả có màu giống nhau.
\choice
{\True $180$}
{$150$}
{$120$}
{$60$}
\loigiai{
\begin{itemize}
\item Số cách lấy được ba quả cầu màu xanh là $3\cdot 4\cdot 5=60$.
\item Số cách lấy được ba quả cầu màu đỏ là $4\cdot 3\cdot 5=60$.
\item Số cách lấy được ba quả cầu màu trắng là $5\cdot 6\cdot 2=60$.
\end{itemize}
Vậy số cách lấy được ba quả cầu cùng màu là $60+60+60=180$.
}
\end{ex}

\begin{ex}%[0D8H2-3]
Có bao nhiêu số chẵn mà mỗi số có $ 4$ chữ số đôi một khác nhau?
\choice
{$ 2520$}
{$ 50000$}
{$ 4500$}
{\True $ 2296$}
\loigiai{
\begin{itemize}
\item Số có $ 4$ chữ số khác nhau đôi một: $ 9\cdot \mathrm{A}_9^3$.
\item Số có $ 4$ chữ số lẻ khác nhau đôi một: $ 5\cdot8\cdot\mathrm{A}_8^2$.
\end{itemize}
Vậy số có $ 4$ chữ số chẵn khác nhau đôi một: $ 9\cdot\mathrm{A}_9^3-5\cdot8\cdot\mathrm{A}_8^2=2296$.}
\end{ex}

\begin{ex}%[Mức độ 2]%[0D8H2-4]
Có $6$ học sinh và $3$ thầy giáo A, B, C. Hỏi có bao nhiêu cách xếp chỗ  $9$ người đó ngồi trên một hàng ngang có $9$ chỗ sao cho mỗi thầy giáo ngồi giữa hai học sinh.
\choice
{$4\,320$}
{$90$}
{\True $43\,200$}
{$720$}
\loigiai{
Sắp $6$ học sinh thành một hàng ngang, giữa $6$ học sinh có $5$ khoảng trống, ta chọn $3$ khoảng trống và mời $3$ giáo viên ngồi vào thì được cách sắp thỏa yêu cầu bài toán.\\
Vậy có $6!\cdot \mathrm{A}_5^3=43\,200$ cách.
}
\end{ex}

\begin{ex}%[0D8H2-6]
Cho tứ diện $ABCD$. Hỏi có bao nhiêu vectơ khác vectơ $\overrightarrow{0}$ mà mỗi vectơ có điểm đầu, điểm cuối là hai đỉnh của tứ diện $ABCD$?
\choice
{\True $12$}
{$4$}
{$10$}
{$8$}
\loigiai{
Ta có mỗi vectơ khác vectơ $\overrightarrow{0}$ có điểm đầu và điểm cuối là $2$ đỉnh trong số $4$ đỉnh của tứ diện là một chỉnh hợp chập $2$ của $4$ phần tử.\\
Vậy số vectơ khác vectơ $\overrightarrow{0}$ có điểm đầu, điểm cuối là hai đỉnh của tứ diện $ABCD$ là $\mathrm{A}_4^2=12$.}
\end{ex}

\begin{ex}%[0D8H3-2]
Nếu khai triển nhị thức Niu-tơn $\left(x-1\right)^5=a_5x^5+a_4x^4+a_3x^3+a_2x^2+a_1x+a_0$
thì tổng $a_5+a_4+a_3+a_2+a_1+a_0$ bằng:
\choice
{$-32$}
{\True $0$}
{$1$}
{$32$}
\loigiai{
Ta có
\begin{align*}
\left(x-1\right)^5&=\mathrm{C}_5^0\cdot (x)^5+\mathrm{C}_5^1(x)^4(-1)+\mathrm{C}_5^2(x)^3(-1)^2+\cdots+\mathrm{C}_5^5(x)^0(-1)^5 \\
& =\mathrm{C}_5^0\cdot x^5-\mathrm{C}_5^1\cdot x^4+\mathrm{C}_5^2\cdot x^3-\mathrm{C}_5^3\cdot x^2+\mathrm{C}_5^4\cdot x^1-\mathrm{C}_5^5\cdot x^0
\end{align*}
Khi đó tổng $a_5+a_4+a_3+a_2+a_1+a_0$ bằng $\mathrm{C}_5^0-\mathrm{C}_5^1+\mathrm{C}_5^2-\mathrm{C}_5^3+\mathrm{C}_5^4-\mathrm{C}_5^5=0$.}
\end{ex}

\begin{ex}%[10-11-12EX-HK1-2425]%[Trương Quan Kía]%[0H9H1-2]
Cho hai véc-tơ $\overrightarrow{a} = (3; -1)$, $\overrightarrow{b} = (1; -2)$. Tọa độ véc-tơ $\overrightarrow{u} = 2\overrightarrow{a} - 3\overrightarrow{b}$ là
\choice
{\True $\overrightarrow{u} = (3; 4)$}
{$\overrightarrow{u} = (-3; 4)$}
{$\overrightarrow{u} = (4; 3)$}
{$\overrightarrow{u} = (5; 3)$}
\loigiai
{
Ta có:
$2\overrightarrow{a} = (6; -2), \quad 3\overrightarrow{b} = (3; -6)$.\\
Vậy $\overrightarrow{u} = (6 - 3; -2 - (-6)) = (3; 4)$.
}
\end{ex}

\begin{ex}%[0H9H3-2]
Cho đường thẳng $d$ có phương trình tham số $\heva{
& x=5+t \\
& y=-9-2t
}$. Phương trình tổng quát của đường thẳng $d$ là
\choice
{\True $2x+y-1=0$}
{$-2x+y-1=0$}
{$x+2y+1=0$}
{$2x+3y-1=0$}
\loigiai{
Đường thẳng $d \colon \heva{
& x=5+t \\
& y=-9-2t \\
}$ $\Leftrightarrow \heva{
& t=x-5 \\
& y=-9-2t \\
}$ $\Rightarrow y=-9-2\left( x-5 \right)$ $\Leftrightarrow 2x+y-1=0$.
}
\end{ex}

\begin{ex}%[0H9H4-2]%[TeX Hoá Đề CK2 - Form 2025 - Đợt 2 - GV: Võ Thị Thùy Trang]
Cho $2$ điểm $A(1;1)$, $B(7;5)$. Phương trình đường tròn đường kính $AB$ là
\choice
{$x^2+y^2+8x+6y-12=0$}
{\True $x^2+y^2-8x-6y+12=0$}
{$x^2+y^2+8x+6y+12=0$}
{$x^2+y^2-8x-6y-12=0$}
\loigiai{
Đường tròn có tâm $I$ là trung điểm của đoạn thẳng $AB$, bán kính $R=\dfrac{AB}{2}$.\\
Suy ra $\heva{&x_I=\dfrac{x_A+x_B}{2}\\&y_I=\dfrac{y_A+y_B}{2}}\Rightarrow \heva{&x_I=\dfrac{1+7}{2}=4\\&y_I=\dfrac{1+5}{2}=3}\Rightarrow I=(4;3)$.\\
$R=\dfrac{AB}{2}=\dfrac{\sqrt{(7-1)^2+(5-1)^2}}{2}=\sqrt{13}$.\\
Phương trình đường tròn đường kính $AB$ là $(x-4)^2+(y-3)^2=(\sqrt{13})^2$\\
$\Leftrightarrow x^2+y^2-8x-6y+12=0$.\\
Kết luận phương trình đường tròn đường kính $AB$ là $x^2+y^2-8x-6y+12=0$.}
\end{ex}

\begin{ex}%[0H9H5-2]%[Dự án đề kiểm tra Toán 11 GHKII NH23-24- Hieu Hieu Minh Minh]%[THPT Toàn Thắng- Hải Phòng]
Trong mặt phẳng tọa độ $O x y$, phương trình nào sau đây là phương trình chính tắc của elip?
\choice
{\True $\dfrac{x^2}{144}+\dfrac {y^2}{25}=1$}
{$\dfrac{x^2}{144}-\dfrac {y^2}{25}=1$}
{$\dfrac{x^2}{25}+\dfrac {y^2}{144}=1$}
{$\dfrac{x^2}{144}+\dfrac {y^2}{25}=-1$}
\loigiai{
Phương trình chính tắc của elip có dạng $\dfrac{x^2}{a^2}+\dfrac{y^2} {b^2}=1$ với $a>b$.\\
Do đó phương trình chính tắc của elip là $\dfrac{x^2}{144}+\dfrac {y^2}{25}=1$.
}
\end{ex}

\begin{ex}%[0H9H5-4]%[Dự án đề kiểm tra Toán khối 10 HK2-NH23-24-Đợt 15-Dương Công Tạo]%[THPT Đặng Huy Trứ - Huế]
Cho hypebol $(H)\colon \dfrac{x^2}{100}-\dfrac{y^2}{64}=1$. $M$ là điểm nằm trên $(H)$, $F_1$, $F_2$ là hai tiêu điểm của $(H)$. Khi đó $\left|MF_1-MF_2\right|$ bằng
\choice
{$10$}
{$16$}
{$100$}
{\True $20$}
\loigiai{
Ta có $(H)\colon \dfrac{x^2}{100}-\dfrac{y^2}{64}=1$ $\Rightarrow a^2=100\Leftrightarrow a=10$.\\
Mà $M\in (H)$ nên $\left|MF_1-MF_2\right|=2a=20$.
}
\end{ex}
\TNTF
\begin{ex}%[0D8H2-4]%[Dự án đề kiểm tra Toán Khối 10 CK2 NH23-24-Dot 16- Khắc Thiên]%[THPT Chuyên Hùng Vương - Phú Thọ]
Tổ $I$ của lớp $10A$ gồm có $7$ học sinh gồm $4$ nam và $3$ nữ.
\choiceTF[1t]
{\True Xếp $7$ học sinh của tổ $I$ vào một hàng ngang để chụp ảnh có $7!$ cách}
{Có $\mathrm{C}_7^2$ cách chọn ra một cặp nam nữ của tổ $I$ để tham gia hát song ca}
{\True Lớp trưởng cần chọn ra $3$ học sinh của tổ $I$ để trực nhật lớp, trong đó $1$ bạn quét lớp, $1$ bạn lau bảng, $1$ bạn kê bàn ghế. Số cách chọn là $\mathrm{A}_7^3$ cách}
{\True Có $720$ cách xếp $7$ học sinh của tổ $I$ vào một hàng dọc sao cho $3$ bạn nữ luôn đứng cạnh nhau}
\loigiai{
\begin{itemchoice}
\itemch Tổ $I$ của lớp $10A$ gồm có $7$ học sinh gồm $4$ nam và $3$ nữ.\\ Xếp $7$ học sinh của tổ $I$ vào một hàng ngang để chụp ảnh có $7!$ (cách).
\itemch Chọn $1$ bạn nam có $4$ cách, chọn $1$ bạn nữ có $3$ cách. \\Theo quy tắc nhân có $4\cdot3=12$ (cách).
\itemch Chọn $1$ bạn quét lớp có $7$ cách, chọn $1$ bạn lau bảng có $6$ cách, chọn $1$ bạn kê bàn ghế $5$ cách.\\Theo quy tắc nhân ta có $7\cdot6\cdot5=\mathrm{A}_{7}^{3}$ (cách).
\itemch Số cách xếp $7$ học sinh của tổ $I$ vào một hàng dọc sao cho $3$ bạn nữ luôn đứng cạnh nhau là $5!\cdot 3!=720$ (cách).
\end{itemchoice}
}
\end{ex}

\begin{ex}%[Dự án 18 - TeamTeXHoa - Phạm Quốc Toàn]%[0H9H4-5]
Trong mặt phẳng $Oxy$, cho đường tròn $\left(C \right) \colon \left(x+2 \right)^2+ \left(y-1 \right)^2=9$ và hai điểm $A \left(-4;3 \right)$, $B \left(2;-1 \right)$. Các mệnh đề sau đúng hay sai?
\choiceTF
{Điểm $A$ nằm trên đường tròn $\left(C \right)$}
{\True Điểm $B$ nằm ngoài đường tròn $\left(C \right)$}
{Phương trình đường thẳng $d$ đi qua điểm $A$ sao cho khoảng cách từ tâm đường tròn đến đường thẳng $d$ là lớn nhất là $x-y-1=0$}
{\True Giá trị lớn nhất của $BM$ với $M$ là điểm chuyển động trên đường tròn là $2\sqrt{5}+3$}
\loigiai{
\begin{itemchoice}
\itemch Ta có $\left(-4+2 \right)^2+ \left(3-1 \right)^2=8 \ne 9$ nên điểm $A$ không nằm trên đường tròn $\left(C \right)$.
\itemch Đường tròn $\left(C \right) \colon \left(x+2 \right)^2+ \left(y-1 \right)^2=9$ có tâm $I \left(-2;1 \right)$ và bán kính $R=3$. \\
Mà $IB=\left| \overrightarrow{IB}\right|=2\sqrt{5}>3$ nên điểm $B$ nằm ngoài đường tròn $\left(C \right)$.
\itemch Gọi $H$ là hình chiếu vuông góc của $I$ lên đường thẳng $d$. Khi đó $IH \le IA$. \\
Do đó khoảng cách lớn nhất là $IA$, khi đó đường thẳng $d$ đi qua $A$ và vuông góc với $IA$. Ta có $\overrightarrow{IA}= \left(-2;2 \right)$. \\
Phương trình đường thẳng $d \colon -2 \left(x+4 \right)+2 \left(y-3 \right)=0 \Leftrightarrow x-y+7=0$.
\itemch Gọi $G$, $K$ lần lượt là giao điểm của đường thẳng $IB$ và đường tròn $\left(C \right)$.
\begin{center}
\begin{tikzpicture}[>=stealth, x=1.0 cm, y=1.0 cm]
\draw[->] (-6,0) --(0,0) node[below left]{$O$} -- (3,0) node[below]{$x$};
\draw[->](0,-3)--(0, 4.5) node[right]{$y$};
\draw(-2,1) circle (3); %(0,0) là tọa độ tâm, 1 là bán kính
\draw(-4.67, 2.28) -- (-2, 1) -- (0, 0) -- (0.68, -0.34) -- (2, -1) (-2, 1) -- (-1, 3.83) -- (2, -1);
\draw (-2,1) circle node[below]{$I$};
\draw (-1,3.83) circle node[above]{$M$};
\draw (-4.68, 2.34) circle node[above left]{$K$};
\draw (0.68, - 0.34) circle node[below]{$G$};
\draw (2, -1) circle node[below]{$B$};
\end{tikzpicture}
\end{center}
Với ba điểm $I$, $B$, $M$ ta có \\
$IB-IM\le MB\le IB+IM\Rightarrow IB-IG\le MB\le IB+IK\Rightarrow BG\le BM\le BK$. \\
Trong đó $BG=IB-IG=2\sqrt{5}-3$, $BK=IB+IK=2\sqrt{5}+3$. \\
Vậy $BM$ đạt giá trị lớn nhất là $2\sqrt{5}+3$ khi $M$ trùng $K$.
\end{itemchoice}
}
\end{ex}
\TNSA
\begin{ex}%[Chuẩn hóa M25, Nguyễn Văn Vũ]%[0D0V2-5]
Cho tập hợp $X=\{0;1;2;3;4;5;6\}$. Từ các chữ số trong tập $X$ có thể lập được bao nhiêu số tự nhiên thỏa mãn có $5$ chữ số đôi một khác nhau.
\shortans{$2160$}
\loigiai{
\item Gọi số cần tìm là $\overline{abcde}$, $a\neq 0$.
\begin{itemize}
\item Chọn $a$ có $6$ cách.
\item Chọn $b,c,d,e$ có $\mathrm{A}_6^4$.\\
Vậy có $6\cdot \mathrm{A}_6^4=2160.$
\end{itemize}
}
\end{ex}

\begin{ex}%[0D8H2-5]%[Dự án Toán thực tế]%[Vũ Hồng Toàn]
Mỗi máy tính tham gia vào mạng phải có một địa chỉ duy nhất, gọi là địa chỉ IP, nhằm định danh máy tính đó trên Internet. Xét tập hợp $A$ gồm các địa chỉ IP có dạng $192$.$168$$ abc$.$deg$, trong đó $a$, $d$ là các chữ số khác nhau được chọn ra từ các chữ số $1$, $2$ còn $b$, $c$, $e$, $g$ là các chữ số đôi một khác nhau được chọn ra từ các chữ số $0,1,2,3,4,5$. Hỏi tập hợp $A$ có bao nhiêu phần tử?
\shortans{$1\,800$}
\loigiai{
$a$, $d$ là các chữ số khác nhau được chọn ra từ các chữ số $1$, $2$ nên có $2$ cách chọn.\\
$b$ và $c$ là các chữ số đôi một khác nhau được chọn ra từ các chữ số $0,1,2,3,4,5$ nên có $\mathrm{A}_6^2$ cách chọn.\\
$e$, $g$ là các chữ số đôi một khác nhau được chọn ra từ các chữ số $0,1,2,3,4,5$ nên có $\mathrm{A}_6^2$ cách chọn.\\
Vậy tập hợp $A$ có số phần tử là $2\cdot \mathrm{A}_6^2\cdot 1\cdot \mathrm{A}_6^2=1\,800$ (phần tử).
}
\end{ex}

\begin{ex}%[0D8H3-4]%[Dự án đề kiểm tra Toán 10, 11 HKII NH23-24-BCTuan]%[THPT LTV-Hà Nội]
Tìm hệ số của số hạng chứa $x^2y^3$ trong khai triển $(2x+y)^5$.
\shortans[3]{$40$}
\loigiai{
Ta có $(2x+y)^5=\displaystyle\sum\limits_{k=0}^5 \mathrm{C}_5^k2^{5-k}x^{5-k}y^k$.\\
Số hạng chứa $x^2y^3$ tương ứng với $k=3$.\\
Hệ số của số hạng là $\mathrm{C}_5^2 2^2=40$.
}
\end{ex}

\begin{ex}%[0H9V4-7]%[Dự án  Toán thực tế]%[Huỳnh Xuân Tín]
\immini[thm]{Hình mô phỏng một trạm thu phát sóng điện thoại di động đặt ở vị trí $I$ có tọa độ $(-2;1)$ trong mặt phẳng toạ độ (đơn vị trên hai trục là ki-lô-mét). Bán kính phủ sóng là $3$ km. Tính theo đường chim bay, xác định khoảng cách ngắn nhất để một người ở vị trí có tọa độ $(-3;4)$ di chuyển được tới vùng phủ sóng theo đơn vị ki-lô-mét (làm tròn kết quả đến hàng phần mười).
}{
\begin{tikzpicture}[line join=round,font=\footnotesize,scale=.6,>=stealth]
\path (0,0) coordinate (O) (-3,4) coordinate (M) (-2,1) coordinate (I)
node[shift={(120:1.2)},text width=2cm,align=center]{\scriptsize Trạm phát sóng};
\draw[-stealth] (-5.5,0)--(2.5,0) node[below] {$x$};
\draw[-stealth] (0,-2)--(0,4.5) node[left] {$y$};
\draw (I) circle (3);
\draw[dashed] (-2,0)node[below]{$-2$}|-(0,1) node[right]{$1$};
\foreach \i/\j in {I/135,O/-135,M/90}
\fill (\i) circle (1.3pt) node[shift={(\j:3mm)}]{$\i$};
\end{tikzpicture}\vspace{2mm}
}
\shortans{$0{,}2$}
\loigiai{
Ta có $IM=\sqrt{(-3+2)^2+(4-1)^2}=\sqrt{10}$ km.\\
Vì $IM>R$ nên $M$ đang ở ngoài vùng phủ sóng.\\
Do đó người ở vị trí $M$ cần di chuyển ít nhất một khoảng bằng $IM-R=\sqrt{10}-3 \approx 0{,}2$ km.
}
\end{ex}
\TL
\begin{ex}%[10-11-GK2-2425]%[Đỗ Minh Phúc]%[0D8H3-2]
Khai triển biểu thức $\left(2 x-\dfrac{3}{2}\right)^5$ bằng nhị thức Newton.
\loigiai
{
Ta có
\begin{eqnarray*}
\left(2x-\dfrac{3}{2}\right)^5&=&\mathrm{C}_{5}^{0}(2x)^5+\mathrm{C}_{5}^{1}(2x)^4\left(-\dfrac{3}{2}\right)+\mathrm{C}_{5}^{2}(2x)^3\left(-\dfrac{3}{2}\right)^2+\mathrm{C}_{5}^{3}(2x)^2\left(-\dfrac{3}{2}\right)^3+\mathrm{C}_{5}^{4}(2x)\left(-\dfrac{3}{2}\right)^4+\mathrm{C}_{5}^{5}\left(-\dfrac{3}{2}\right)^5\\
&=&32x^5-120x^4+180x^3-135x^2+\dfrac{405}{8}x-\dfrac{243}{32}.
\end{eqnarray*}

}
\end{ex}

\begin{ex}%[0D8V2-4]%[Dự án đề kiểm tra Toán 10 HKII NH23-24- Hiếu Phan]%[THPT Phú Nhuận]
Hội đồng của một công ty TNHH gồm $16$ người trong đó có $6$ nữ, từ hội đồng này người ta bầu ra $1$ chủ tịch, $1$ phó chủ tịch và $3$ ủy viên. Hỏi có bao nhiêu cách bầu chọn sao cho trong $5$ người được bầu nhất thiết phải có nữ?
\loigiai
{
Ta có $n(\Omega)=\mathrm{C}_{16}^5$.\\
Số cách chọn $5$ người không có nữ là $\mathrm{C}_{10}^5$.\\
Suy ra có $\mathrm{C}_{16}^5-\mathrm{C}_{10}^5=4116$ cách chọn sao cho trong $5$ người được bầu nhất thiết phải có nữ.
}
\end{ex}

\begin{ex}%[1 điểm]%[0H9V4-5]%[Dự án đề kiểm tra Toán 10 HKII NH23-24- Anh Hoàng Lê]%[THPT Chuyên Lê Hồng Phong - Tp HCM]
Trong mặt phẳng $Oxy$, cho đường tròn $(C)\colon x^2+y^2-2x+4y-4=0$ và đường thẳng $d\colon 2x+y+2 \sqrt{15}=0$. Có bao nhiêu điểm thuộc $d$ sao cho qua điểm đó có thể kẻ hai tiếp tuyến đến $(C)$ và góc giữa hai tiếp tuyến bằng $60^{\circ}$.
\loigiai{
Gọi $M$ là điểm thỏa yêu cầu bài toán. Cho tiếp tuyến tại $M$ tiếp xúc với $(C)$ tại $A$ và $B$. Ta chia thành hai trường hợp.

\faStar~~\textbf{Trường hợp 1.} $\widehat{AMB}=60^{\circ}$.\\
Ta có $(C)\colon x^2+y^2-2x+4y-4=0$ hay $(x-1)^2+(y+2)^2=9$ suy ra $(C)$ có tâm là $I(1;-2)$ và bán kính $R=3$.\\
Do $M \in d$ nên ta gọi tọa độ của điểm $M$ là $\left(k;-2k-2 \sqrt{15}\right)$ suy ra
\[I M=|\overrightarrow{IM}|=\sqrt{\left(k-1\right)^2+\left(-2k-2 \sqrt{15}+2\right)^2}. \]
Mặt khác  $\widehat{IMA}=\widehat{IMB}=\dfrac{\widehat{AMB}}{2}=\dfrac{60^{\circ}}{2}=30^{\circ}$ và $I A \perp A M$ nên
\[I M=\dfrac{IA}{\sin \widehat{AMI}}=\dfrac{R}{\sin 30^{\circ}}= \dfrac{3}{\sin 30^{\circ}}=6.\] Suy ra
\[
\left(k-1\right)^2+\left(-2k-2 \sqrt{15}+2\right)^2=36 \,\, \text{hay} \,\,  5  k^2-2k\left(5-4 \sqrt{15}\right)+29-8 \sqrt{15}=0.
\]
Ta có $\Delta'=\left(5-4 \sqrt{15}\right)^2-5 \cdot \left(29-8 \sqrt{15}\right)=120>0$.\\
Vậy phương trình trên có hai nghiệm phân biệt hay có  $2$  điểm $M$ thỏa.

\faStar~~\textbf{Trường hợp 2.} $\widehat{AMB}=120^{\circ}$.\\
Tương tự như trên ta có $I M=\dfrac{3}{\sin 60^{\circ}}=2 \sqrt{3}$ nên ta có phương trình
\[
\left(k-1\right)^2+\left(-2k-2 \sqrt{15}+2\right)^2=12 \,\, \text{hay} \,\,  5k^2-2k\left(5-4 \sqrt{15}\right)+53-8 \sqrt{15}=0.
\]
Ta có $\Delta'=\left(5-4 \sqrt{15}\right)^2-5 \cdot \left(53-8 \sqrt{15}\right)=0$ hay phương trình trên có nghiệm kép.\\
Vậy có tất cả  $3$  điểm $M$ thỏa đề bài.
}
\end{ex}

